\documentclass[11pt]{article}
\usepackage{graphicx}
\usepackage{amsfonts}
\usepackage{amsmath}
\setlength{\topmargin}{0.0in}     % top of paper to head (less one inch)
\setlength{\headheight}{0in}    % height of the head
\setlength{\headsep}{0in}       % head to the top of the body
%\setlength{\textheight}{9.5in} % height of the body
\setlength{\textheight}{9in} % height of the body
\setlength{\oddsidemargin}{0mm} % left edge of paper to body (less one inch)
\setlength{\evensidemargin}{0mm} % ditto, even pages
\setlength{\textwidth}{6.5in}   % width of body
\setlength{\topskip}{0in}       % top of body to bottom of first line of text
%\setlength{\footheight}{0.25in} % height of footer
\setlength{\footskip}{0.50in}   % bottom of text to bottom of foot
% \renewcommand{\baselinestretch}{1.4}   % uncomment to double space
\setlength{\parskip}{1ex}
\setlength{\parindent}{0in}


\newtheorem{theorem}{Theorem}
\newtheorem{corollary}{Corollary}
\newtheorem{lemma}{Lemma}
\newtheorem{observation}{Observation}
\newtheorem{definition}{Definition}
\newtheorem{fact}{Fact}
\newcommand{\proof}{\vspace*{-1ex} \noindent {\bf Proof: }}
\newcommand{\proofsketch}{\vspace*{-1ex} \noindent {\bf Proof Sketch: }}
\newcommand{\qed}{\hfill\rule{2mm}{2mm}}
\newcommand{\ceiling}[1]{{\left\lceil{#1}\right\rceil}}
\newcommand{\floor}[1]{{\left\lfloor{#1}\right\rfloor}}
\newcommand{\paren}[1]{\left({#1}\right)}
\newcommand{\braces}[1]{\left\{{#1}\right\}}
\newcommand{\brackets}[1]{\left[{#1}\right]}
\newcommand{\Prob}{{\rm Prob}}
\newcommand{\prob}{{\rm Prob}}
%
\newcommand{\ob}{$\langle$}
\newcommand{\cb}{$\rangle$}
\newcommand{\mymod}{\mbox{ mod }}

\newcommand{\header}[3]{
   \pagestyle{plain}
   \noindent
   \begin{center}
   \framebox{
      \vbox{
    \hbox to 6.28in { {\bf CS 334 Computer Security
\hfill Fall 2014} }
       \vspace{4mm}
       \hbox to 6.28in { {\Large \hfill #1 \hfill} }
       \vspace{4mm}
    \hbox to 6.28in { {\sl #2 \hfill #3} }
      }
   }

   \end{center}
   \vspace*{4mm}
}

%\input{psfig.tex}

\begin{document}

\header{Common Implementation Flaws --- Memory Safety}{Prof. Szajda}

\textit{Thanks to Vern Paxson and Dave Wagner of UC-Berkeley for providing this
handout.}

In the next few lectures we will be looking at software
security problems associated with software implementation. Even with
a perfect design, perfect specification, and perfect algorithms,
you may still have implementation vulnerabilities. In fact, after
configuration errors, implementation errors are probably the largest
single class of security errors exploited in practice.  The purpose of
the next few lectures is to teach you about software security, and in particular
about the kinds of implementation errors that can lead to serious vulnerabilities.

We will start by showing you some of the most common
implementation flaws. Because many security-critical
applications have been written in C, and because C has peculiar
pitfalls of its own, many of these examples will be
C-specific.  Don't let this fool you into believing that 
other languages are safe.  Some, like Java, take pains to
eliminate certain types of implementation flaws.  No language,
however, is completely safe from these errors.  In addition,
implementation flaws can occur at all levels: in
improper use of the programming language, the libraries, the operating
system, or in the application logic. Since by far the most common class of
implementation flaw is the buffer overrun, we will start there.

\section{Buffer Overruns}

A \textit{buffer overflow} (also called \textit{buffer overrun})
vulnerability is present when code has been written such that it is
possible for the user to place more data into memory (typically a
variable) than the amount of space that has been allocated for that
data by the programmer.  In effect, the programmer has neglected to
check that what is being received fits into the space in which it is
to be placed.  The result is that memory adjacent to the desired
storage space is overwritten.  Buffer overflow vulnerabilities are a
particular risk in C, since C does not, by default, provide bounds
checks on many operations that write to memory.  Since C is an
especially widely used systems programming language, you might not be
surprised to hear that buffer overflows are one of the most pervasive
kind of implementation flaws around.

As a low-level language, we can think of
C as a portable assembly language. The programmer is exposed to the
bare machine (which is one reason that C is such a popular systems
language). As mentioned above, a particular weakness is the absence
of automatic bounds-checking, in particular for array or pointer access. 
If the programmer fails to perform adequate
bounds checks, the user can trigger an out-of-bounds memory access that writes
beyond the bounds of some memory region.  Attackers can use these
out-of-bounds memory accesses to corrupt the program’s intended
behavior.  

Let us start with a simple example.

\begin{ttfamily}
\begin{tabbing}
12345\=first\=second \kill
char buf[80];\\
void vulnerable() \{ \\
\>gets(buf);\\
\}\\
\end{tabbing}
\end{ttfamily}

In this example, \texttt{gets()} reads as many bytes of input as are available
on standard input, and stores them into \texttt{buf[]}. If the input contains
more than 80 bytes of data, then \texttt{gets()} will write past the end of
\texttt{buf}, overwriting some other part of memory. This is a bug. Obviously,
this bug might cause the program to crash or core-dump.  What might be
less obvious is that the consequences can be far worse than that.

To illustrate some of the dangers, we modify
the example slightly.

\begin{ttfamily}
\begin{tabbing}
12345\=first\=second \kill
char buf[80];\\
int authenticated = 0;\\
void vulnerable() \{\\
\>gets(buf);\\
\}\\
\end{tabbing}
\end{ttfamily}

Imagine that elsewhere in the code is a login routine that sets the
\texttt{authenticated} flag only if the user proves knowledge of the
password.  Unfortunately, the \texttt{authenticated} flag is stored in memory
right after \texttt{buf}. If the attacker can write 81 bytes of data to \texttt{buf}
(with the 81st byte set to a non-zero value), then this will set the
\texttt{authenticated} flag to true, and the attacker will gain access. The
program above allows that to happen, because \texttt{gets} does no
bounds-checking: it will write as much data to \texttt{buf} as is supplied to
it. In other words, the code above is vulnerable: an attacker who can
control the input to the program, can bypass the password checks.

Now consider another variation:

\begin{ttfamily}
\begin{tabbing}
12345\=first\=second \kill
char buf[80]; \\
int (*fnptr)(); \\
void vulnerable() \{ \\
\>gets(buf);\\
\}\\
\end{tabbing}
\end{ttfamily}

The function pointer \texttt{fnptr} is invoked elsewhere in the program (not
shown). This enables a more serious attack: the attacker can overwrite
\texttt{fnptr} with any address of her choosing, and redirecting program execution 
to some other memory location. A crafty attacker could supply
an input that consists of malicious machine instructions, followed by
a few bytes that overwrite \texttt{fnptr} with some address $A$.  When 
\texttt{fnptr} is next invoked, the flow of control is redirected to address $A$. Notice
that in this attack, the attacker can choose the address $A$ however she
likes---so, for instance, she can choose to overwrite \texttt{fnptr} with an
address where the malicious machine instructions will be stored (e.g.,
the address \texttt{\&buf[0]}). This is a \textit{malicious code injection} attack. 
Of course, many variations on this attack are possible: for instance, the
attacker could arrange to store the malicious code anywhere else
(e.g., in some other input buffer), rather than in \texttt{buf}, and redirect
execution to that other location.  

Malicious code injection attacks
allow an attacker to sieze control of the program. At the conclusion
of the attack, the program is still running, but now it is executing
code chosen by the attacker, rather than the original software. For
instance, consider a web server that receives requests from clients
across the network and processes them. If the web server contains a
buffer overrun in the code that processes such requests, a malicious
client would be able to seize control of the web server process. If
the web server is running as root, once the attacker seizes control,
the attacker can do anything that root can do; for instance, the
attacker can leave a backdoor that allows her to log in as root
later. At that point, the system has been ``owned.''

Buffer overflow vulnerabilities and malicious code injection are a
favorite method used by worm writers and attackers. A pure example
such as we showed above is relatively rare. However, it does occur:
for instance, it formed the vectors used by the first major Internet
worm (the Morris worm). Morris took advantage of a buffer overflow in
\texttt{in.fingerd} (the network finger daemon) to overwrite the
filename of a command executed by \texttt{in.fingerd}, similar to the
example above involving an overwrite of an ``authenticated'' flag. But
pure attacks, as illustrated above, are only possible when the code
satisfies certain special conditions: the buffer that can be
overflowed must be followed in memory by some security-critical data
(e.g., a function pointer, a flag that has a critical influence on the
subsequent flow of execution of the program). Because these conditions
occur only rarely in practice, attackers have developed more effective
methods of malicious code injection.

\section{Stack smashing}

One powerful method for exploiting buffer overrun vulnerabilities
takes advantage of the way local variables are layed out on the stack.

We need to review some background material first.  Let's recall C's memory
layout:\\
\vspace*{-2.5in}
\begin{center}
\setlength{\unitlength}{1.0in}
\hspace*{-2.5in}\begin{picture}(0.5, 3.0)
%\put(0.1,0.3){text region}
\put(0.0, 0.25){\framebox(1.0,0.25){text region}}
\put(1.0, 0.25){\framebox(2.0,0.25){heap $\qquad \cdots \qquad$ stack}}
\put(0.0, 0.0){\makebox(1.0,0.25){0x00..0}}
\put(1.0, 0.0){\makebox(2.0,0.25){\hspace*{1.0in} 0xFF..F}}
\end{picture} 
\end{center}

The text region contains the executable code of the program. The heap
stores dynamically allocated data (and grows and shrinks as objects
are allocated and freed). Local variables are stored and other
information associated with each function call is stored on the stack
(which grows and shrinks with function calls and returns). In the
picture above, the text region starts at smaller-numbered memory
addresses (e.g., 0x00..0), and the stack region ends at
larger-numbered memory addresses (0xFF..F).

Function calls push new
stack frames onto the stack. A stack frame includes space for all the
local variables used by that function, and other book-keeping
information used by the compiler for this function invocation. On
Intel (x86) machines, the stack grows down. This means that the stack
grows towards smaller memory addresses. There is a special register,
called the stack pointer (SP), that points to the beginning of the
current stack frame. Thus, the stack extends from the address given in
the SP until the end of the memory allocated for the process.
Pushing a new frame on the stack involves subtracting the length of that frame from SP.  

Intel (x86)
machines have another special register, called the instruction pointer
(IP), that points to the next machine instruction to execute. For most
machine instructions, the machine reads the instruction pointed to by
IP, executes that instruction, and then increments the IP. Function
calls cause the IP to be updated differently: the current value of IP
is pushed onto the stack (this will be called the return address), and
the program then jumps to the beginning of the function to be
called. The compiler inserts a function prologue---some automatically
generated code that performs the above operations---into each function,
so it is the first thing to be executed when the function is
called. The prologue pushes the current value of SP onto the stack and
allocates stack space for local variables by decrementing the SP by
some appropriate amount.  

When the function returns, the old SP and
return address are retrieved from the stack, and the stack frame is
popped from the stack (by restoring the old SP value). Execution
continues from the return address.  

After all that background, we're now ready to see how a 
\textit{stack smashing} attack works. Suppose the code looks like this:

\begin{ttfamily}
\begin{tabbing}
12345\=first\=second \kill
void vulnerable() \{ \\
\>char buf[80]; \\
\>gets(buf); \\
\} \\
\end{tabbing}
\end{ttfamily}

When \texttt{vulnerable()} is called, a stack frame is pushed onto the
stack.  The stack will look something like this:
\vspace*{-3.5in}
\begin{center}
\setlength{\unitlength}{1.0in}
\hspace*{-3.0in}\begin{picture}(0.25, 4.0)
%\put(0.1,0.3){text region}
\put(0.0, 0.0){\framebox(2.0,0.25){\texttt{buf} \hspace*{0.2in}
saved SP \hspace*{0.15in}ret addr}}
\put(1.96, 0.0){\framebox(1.54,0.25){caller's stack frame}}
\put(3.5,0.0){\makebox(0.5,0.25){$\cdots$}}
\put(3.5,0.0){\line(1,0){0.5}}
\put(3.5,0.25){\line(1,0){0.5}}
\end{picture} 
\end{center}

If the input is too long, the code will write past the end of
\texttt{buf} and the saved SP and return address will be overwritten.
This is a \textit{stack smashing} attack.

Stack smashing can be used for malicious code injection. First, the
attacker arranges to infiltrate a malicious code sequence somewhere in
the program's address space, at a known address (perhaps using
techniques previously mentioned).  Next, the attacker provides a
carefully-chosen 88-byte input, where the last four bytes hold the
address of the malicious code\footnote{In this example, I am assuming
  a 32-bit architecture, so that SP and the return address are 4 bytes
  long. On a 64-bit architecture, SP and the return address would be 8
  bytes long, and the attacker would need to provide a 96-byte input
  whose last 8 bytes were chosen carefully to contain the address of
  the malicious code.}.
The \texttt{gets()} call will overwrite the
return address on the stack with the last 4 bytes of the input---in
other words, with the address of the malicious code. When \texttt{vulnerable()}
returns, the CPU will retrieve the return address stored on the stack
and transfer control to that address, handing control over to the
attacker's malicious code.  

The discussion above has barely scratched
the surface of techniques for exploiting buffer overrun bugs. Stack
smashing was first well-documented in 1996 (see ``Smashing the Stack for Fun
and Profit'' Aleph One), though the general idea of overflowing a buffer to
control execution has been known since at least 1972. 
Modern methods are considerably more
sophisticated and powerful. These attacks may seem esoteric, but
attackers have become highly skilled at exploiting them. Indeed, you
can find tutorials on the web explaining how to deal with
complications such as:
\begin{itemize}
\item The malicious code is stored at an unknown
location.
\item There is no way to introduce malicious code into the
program's address space.
\item The buffer is stored on the heap instead of
on the stack.
\item The attack can only overflow the buffer by a single
  byte\footnote{In one of the most impressive instances of this, 
    the Apache webserver at one point had an off-by-one bug
    that allowed an attacker to zero out the next byte immediately
    after the end of the buffer. The Apache developers downplayed the
    impact of this bug and assumed it was harmless---until someone came
    up with a clever and sneaky way to mount a malicious code
    injection attack, using only the ability to write a single zero
    byte one byte past the end of the buffer.}.
\item The characters that can be written to the buffer are limited
  (e.g., to only lowercase letters)\footnote{Imagine writing a
    malicious sequence of instructions where every byte in the
    machine code has to be in the range 0x61 to 0x7A ('a' to
    'z'). Yes, it’s been done.}.  
\end{itemize}

Buffer overrun attacks may appear mysterious or complex or hard to
exploit, but in reality, they are none of the above. Attackers exploit
these bugs all the time. For example, the Code Red II worm compromised
250K machines by exploiting a buffer overflow bug in the IIS web
server. In the past, many security researcher have underestimated the
opportunities for obscure and sophisticated attacks, only to later
discover that the ability of attacker to find clever ways to exploit
these bugs exceeded their imaginations. Attacks once thought to be
esoteric to worry about are now considered easy and routinely mounted
by attackers. The bottom line is this: If your program has a buffer
overflow bug, you should assume that the bug is exploitable and an
attacker can take control of your program.  

How do you avoid buffer
overflows in your code? One way is to check that there is sufficient
space for what you will write before performing the write.

\section{Format String Vulnerabilities}

Let's look at another type of vulnerability.

\begin{ttfamily}
\begin{tabbing}
12345\=first\=second \kill
void vulnerable() \{ \\
\>char buf[80]; \\
\>if (fgets(buf, sizeof buf, stdin) == NULL) \\
\>\>return; \\
\>printf(buf); \\
\} \\
\end{tabbing}
\end{ttfamily}

Do you see the bug? The last line should be
\texttt{printf("\%s",buf)}.  Instead, \texttt{buf} was wrongly passed
as the format string to \texttt{printf()}.  Notice that if
\texttt{buf} contains any \texttt{\%} format-specifiers, \texttt{printf()} will
look for arguments that aren't there, and may crash or core-dump the
program when it follows an invalid pointer. But actually, it's worse
than that. As you may have begun to suspect by now, any time that a
program crashes or core-dumps, it is worth looking closely to see
whether there is something more serious possible.  Let's look.

First, we can see that it may be possible for
an attacker to learn information about the contents of the function's
stack frame. By specifying a string like \texttt{"\%x:\%x"}, an attacker who can
see what is printed has the chance to view the first two words of
stack memory.  

In fact, we can refine this a bit. Suppose we supply a
string like \texttt{"\%x:\%x:\%s"}. What does this do? It prints out the first two
words of stack memory, then treats the next word of stack memory as a
memory address and prints out everything at that address until the
first \verb=`\0'= byte. 
Where does that last word of stack memory come from?
Well, it comes from somewhere in the stack frame for \texttt{printf()}, or, if
we supply enough \texttt{\%x} specifiers to walk past the end of 
\texttt{printf()}'s
stack frame, then it comes from somewhere in \texttt{vulnerable()}'s stack
frame.  

Ah-hah! Now we are on to something. Notice that \texttt{buf} is stored
in \texttt{vulnerable()}'s stack frame, and the attacker can control
the contents of \texttt{buf}, so this means the attacker can control
part of the contents of \texttt{vulnerable()}'s stack frame --- but
that's exactly where the \texttt{\%s} specifier is getting its memory
address from. So, the attacker can store a memory address somewhere in
\texttt{buf}, then arrange that when the \texttt{\%s} specifier reads
a word from the stack to get a memory address, it will receive the
memory address he thoughtfully put there for it. The exploit will
involve supplying a string like
\texttt{\mbox{"\\x04\\x03\\x02\\x01:\%x:\%x:\%x:\%x:\%s"}}.  If the
attacker has arranged to have just the right number of \texttt{\%x}
specifiers, then the address that is read will be read from the first
word of \texttt{buf}, which he has arranged to contain the value
\texttt{0x01020304} (it looks like it has been reversed, but that is
because values are stored in little-endian format on x86 machines). We
can see that in this way, the attacker can pick any desired memory
address, and read out the contents of memory starting at that
address. Thus, an attacker can exploit a format string vulnerability
to learn any secrets stored in the victim's address
space --- cryptographic keys, passwords, they're all potential targets.

It gets worse than that. Thanks to an obscure format specifier 
(\texttt{\%n}), it is possible for an attacker to write any desired value to any
desired address in the victim's memory, if the victim has a format
string bug. I won't bother you with the details of how it is done;
suffice it to say that it is possible. This is now enough to allow
attackers to mount malicious code injection attacks. For instance, the
attacker can introduce a malicious code sequence anywhere into the
victim's memory, then use a format string bug to overwrite a return
address on the stack (or a function pointer) so that it now points to
the beginning of the malicious code.  

Consequently, any program that
contains a format string bug can usually be exploited by the attacker
to take control of the victim program and all privileges it has been
granted on the target system. Format string bugs are, like buffer
overruns, nasty business.

The bottom line: if your program has a format string bug, assume that
the attacker can learn all secrets stored in memory, and assume that
the attacker can take control of your program.

\section{Integer Overflow and Implicit Cast Vulnerabilities}

What's wrong with this code?

\begin{ttfamily}
\begin{tabbing}
12345\=first\=second \kill
char buf[80]; \\
void vulnerable() \{ \\
\>int len = read\_int\_from\_network(); \\
\>char *p = read\_string\_from\_network(); \\
\>if (len > sizeof buf) \{ \\
\>\>error("length too large: bad dog, no cookie for you!"); \\
\>\>return; \\
\>\} \\
\>memcpy(buf, p, len);\\
\}\\
\end{tabbing}
\end{ttfamily}

Here's a hint.  The prototype for \texttt{memcpy()} is:

\noindent \texttt{void *memcpy(void *dest, const void *src, size\_t n);}\\

And the definition of \texttt{size\_t} is:

\texttt{typedef unsigned int size\_t;}

Do you see the bug now? If the attacker provides a negative value for
\texttt{len}, the \texttt{if} statement won't notice anything wrong, and 
\texttt{memcpy()} will be executed with a negative third argument.  
C will cast this negative value to an \texttt{unsigned int} and
it will become a very large positive integer.  Thus \texttt{memcpy()} will copy a huge
amount of memory into \texttt{buf}, overflowing the buffer.

Note that the C compiler won't warn about the type mismatch between
\texttt{signed int} and \texttt{unsigned int}; it silently inserts an implicit
cast. This kind of bug can be hard to spot. The above example is
particularly nasty, because on the surface it appears that the
programmer has applied the correct bounds checks, but they are flawed.

Here is another example. What's wrong with this code?

\begin{ttfamily}
\begin{tabbing}
12345\=first\=second \kill
void vulnerable() \{ \\
\>size\_t len; \\
\>char *buf; \\
\\
\>len = read\_int\_from\_network(); \\
\>buf = malloc(len+5); \\
\>read(fd, buf, len); \\
\>... \\
\}\\
\end{tabbing}
\end{ttfamily}

This code seems to avoid buffer overflow problems, (indeed, it
allocates 5 more bytes than necessary).  And there are no signed/unsigned
problems here, since every integer in sight is unsigned.
But, there is a subtle
problem: len+5 can wrap around if len is too large. For instance, if
len = 0xFFFFFFFF, then the value of len+5 is 4 (on 32-bit
platforms). In this case, the code allocates a 4-byte buffer and then
writes a lot more than 4 bytes into it: a classic buffer overflow. You
have to know the semantics of your programming language very well to
avoid all the pitfalls.

\section{Memory Safety}

Buffer overflow, format string, and the other examples above are
examples of \textit{memory safety} bugs: cases where an attacker can read or
write beyond the valid range of memory regions. Other examples of
memory safety violations include using a dangling pointer (a pointer
into a memory region that has been freed and is no longer valid) and
double-free bugs (where a dynamically allocated object is explicitly
freed multiple times).  C and C++ rely upon the programmer to preserve
memory safety, but bugs in the code can lead to violations of memory
safety.  History has taught us that memory safety violations often
enable malicious code injection and other kinds of attacks.  

Some modern languages are designed to be intrinsically memory-safe, no
matter what the programmer does. Java is one example. Thus,
memory-safe languages eliminate the opportunity for one kind of
programming mistake that has been known to cause serious security
problems.  

We've only scratched the surface of implementation
vulnerabilities. If this makes you a bit more cautious when you write
code, then good! In future lectures we'll discuss how to prevent (or
reduce the likelihood) of these types of flaws and to improve the odds
of surviving any flaws that do creep in.

\end{document}









